\chapter*{Útdráttur}
\addcontentsline{toc}{chapter}{Útdráttur}
\draftchange{Í þessari ritgerð er fjallað um hönnun, smíði og mat á orkunýtnu þráðlausu skynjaraneti sem byggir á LoRa-fjarskiptum og er ætlað til að fylgjast með jarðvegs- og umhverfisaðstæðum í íslenskum torfþökum. Kerfið samanstendur af skynjaraeiningum sem byggja á Arduino Pro Mini-stýrieiningum, með rýmdarmælum fyrir raka í jarðvegi, DS18B20-hitanemum fyrir jarðvegshita, AHT10-nemum fyrir lofthita og rakastig, staðbundinni gagnageymslu á SD-korti og 433 MHz LoRa-fjarskiptum. Móttökustöð sendir mælingarnar áfram til ThingSpeak-skýjaþjónustunnar, þar sem hægt er að nálgast og skoða gögnin.}

\draftchange{Prófanir í rannsóknarstofu og prófanir við raunverulegar aðstæður voru framkvæmdar til að meta orkunotkun, fjarskiptaafköst, kvörðun skynjara og gagnaupphleðslu í skýið. Við skilgreindan rekstrarferil var meðalstraumnotkun skynjaraeiningar um 3,08 mA. Með þremur 3.500 mAh rafhlöðufrumum tengdum samhliða, með samtals uppgefinni afkastagetu upp á 10.500 mAh, var áætlaður rekstrartími um 142 dagar. Þetta er fræðilegt mat sem tekur ekki tillit til minnkunar á nýtanlegri rafhlöðugetu við lágt hitastig.}

\draftchange{Við prófanir í hálfþéttbýlu umhverfi í Völlum í Hafnarfirði mældust vel heppnuð LoRa-fjarskipti yfir eina tengingu á nokkrum prófunarstöðum innan við 3 km fjarlægða með 433 MHz loftnetum. Mæld RSSI-gildi voru á bilinu $-103$ dBm til $-113$ dBm. Afköst fjarskiptanna voru háð staðsetningu loftneta og hindrunum í umhverfinu og stöku pakkatap sást á svæðum með meiri hindrunum. Við kvörðun rýmdarmælis fyrir jarðvegsraka fékkst annars stigs samband á prófuðu rakabili, 0--40\%, með um það bil 3,01 prósentustiga kvaðratrótarmeðalskekkju (RMSE) miðað við þyngdarmælingar á sömu kvörðunargögnum.}

\draftchange{Upphleðslur gagna til ThingSpeak tóku að meðaltali um 1--2 sekúndur. Kerfið var einnig sett upp á Keldum á Rangárvöllum til að fylgjast með umhverfisaðstæðum í fjórum torfþökum. Gildar mælingar spönnuðu frá fyrstu til síðustu skráningu um það bil 51 til 130 daga, eftir skynjaraeiningum. Niðurstöðurnar sýna að raunhæft er að þróa ódýrt vöktunarkerfi fyrir íslensk torfþök, en jafnframt þarf að taka mið af umhverfisaðstæðum og áreiðanleika vélbúnaðar til lengri tíma.}

\vfill
\textbf{Lykilorð:} \draftchange{þráðlaust skynjaranet, LoRa, eins tengingar fjarskipti, jarðvegsraki, torfþök, umhverfisvöktun, orkunýting, ThingSpeak}

